\subsubsection{Unitary and Orthogonal Structure}

Orthonormal bases allow geometric structure to coexist with coordinate computations. They support triangular forms, and they also turn the dual action of a linear map into its adjoint.

\begin{corollary} \label{cor:upper-triangular-orthonormal-basis} \hlt{Upper-Triangular Form in an Orthonormal Basis}\\
Let $\mathcal{V}$ be a finite-dimensional inner-product space and let $T\in\LL(\mathcal{V})$. If $T$ has an upper-triangular matrix with respect to some basis of $\mathcal{V}$, then $T$ has an upper-triangular matrix with respect to some orthonormal basis of $\mathcal{V}$.
\begin{proof}
Let $\beta=(\mathbf{v}_1,\ldots,\mathbf{v}_n)$ be a basis for which $[T]_{\beta\leftarrow\beta}$ is upper triangular. By Theorem \ref{thm:upper-triangular-criteria}, each subspace
\begin{align}
\mathcal{U}_j=\operatorname{span}(\mathbf{v}_1,\ldots,\mathbf{v}_j),\qquad j=1,\ldots,n,
\nonumber
\end{align}
is invariant under $T$. Apply Gram--Schmidt to $\beta$ and obtain an orthonormal basis $\gamma=(\mathbf{e}_1,\ldots,\mathbf{e}_n)$ with
\begin{align}
\operatorname{span}(\mathbf{e}_1,\ldots,\mathbf{e}_j)=\operatorname{span}(\mathbf{v}_1,\ldots,\mathbf{v}_j)=\mathcal{U}_j.
\nonumber
\end{align}
Thus each $\operatorname{span}(\mathbf{e}_1,\ldots,\mathbf{e}_j)$ is invariant under $T$, and Theorem \ref{thm:upper-triangular-criteria} shows that $[T]_{\gamma\leftarrow\gamma}$ is upper triangular.\\
\hlt{Structural Commentary}: Upper triangularity is preservation of a nested flag of subspaces. Gram--Schmidt preserves that flag while making the basis orthonormal.\\
\hlt{Key Insight}: Triangular form depends on the flag, not on the particular basis used inside the flag.
\end{proof}
\end{corollary}

\begin{theorem} \label{thm:schur-triangularization} \hlt{Schur Triangularisation}\\
Let $\mathcal{V}$ be a finite-dimensional complex inner-product space and let $T\in\LL(\mathcal{V})$. Then $T$ has an upper-triangular matrix with respect to some orthonormal basis of $\mathcal{V}$.
\begin{proof}
By Theorem \ref{thm:complex-upper-triangularization}, $T$ has an upper-triangular matrix with respect to some basis of $\mathcal{V}$. Corollary \ref{cor:upper-triangular-orthonormal-basis} replaces that basis by an orthonormal basis without losing upper triangularity.\\
\hlt{Structural Commentary}: Complex triangularisation gives a flag; Gram--Schmidt orthonormalizes it.\\
\hlt{Key Insight}: Schur triangularisation is ordinary triangularisation plus orthonormalization of the flag.
\end{proof}
\end{theorem}

\begin{definition} \label{def:adjoint-linear-map} \hlt{Adjoint of a Linear Map}\\
Let $\mathcal{V}$ and $\mathcal{W}$ be inner-product spaces and let $T\in\LL(\mathcal{V},\mathcal{W})$.\\
A function $S:\mathcal{W}\to\mathcal{V}$ is an \hlt{adjoint} of $T$ if $\innerprod[\mathcal{W}]{T\mathbf{v},\mathbf{w}} =\innerprod[\mathcal{V}]{\mathbf{v},S\mathbf{w}}$, $\mathbf{v}\in\mathcal{V},\ \mathbf{w}\in\mathcal{W}$.
\end{definition}

\begin{lemma} \label{lem:adjoint-unique-if-exists} \hlt{Uniqueness of Adjoints}\\
A linear map has at most one adjoint.
\begin{proof}
Suppose $S_1,S_2:\mathcal{W}\to\mathcal{V}$ are adjoints of $T$. For each $\mathbf{w}\in\mathcal{W}$ and every $\mathbf{v}\in\mathcal{V}$, $\innerprod[\mathcal{V}]{\mathbf{v},S_1\mathbf{w}-S_2\mathbf{w}}=0$.\\
Taking $\mathbf{v}=S_1\mathbf{w}-S_2\mathbf{w}$ gives $S_1\mathbf{w}=S_2\mathbf{w}$. Hence $S_1=S_2$.\\
\hlt{Structural Commentary}: Same tests force candidate adjoints to agree.\\
\hlt{Key Insight}: Inner products separate vectors.
\end{proof}
\end{lemma}

\begin{theorem} \label{thm:finite-dimensional-adjoint-exists} \hlt{Existence of the Adjoint}\\
Let $\mathcal{V}$ be a finite-dimensional inner-product space, let $\mathcal{W}$ be an inner-product space, and let $T\in\LL(\mathcal{V},\mathcal{W})$. Then $T$ has a unique adjoint $T^*:\mathcal{W}\to\mathcal{V}$, and $T^*$ is linear.
\begin{proof}
For each $\mathbf{w}\in\mathcal{W}$, the map $\varphi_{\mathbf{w}}:\mathcal{V}\to\mathbb{F}$ defined by $\varphi_{\mathbf{w}}(\mathbf{v})=\innerprod[\mathcal{W}]{T\mathbf{v},\mathbf{w}}$ is a linear functional. By Theorem \ref{thm:riesz-representation-finite-dimensional}, there is a unique vector $T^*\mathbf{w}\in\mathcal{V}$ such that $\innerprod[\mathcal{W}]{T\mathbf{v},\mathbf{w}} =\innerprod[\mathcal{V}]{\mathbf{v},T^*\mathbf{w}}$, $\mathbf{v}\in\mathcal{V}$.\\
Thus $T^*$ is an adjoint, and it is unique by Lemma \ref{lem:adjoint-unique-if-exists}. For $\mathbf{w}_1,\mathbf{w}_2\in\mathcal{W}$ and $a\in\mathbb{F}$,
\begin{align}
\innerprod[\mathcal{V}]{\mathbf{v},T^*(a\mathbf{w}_1+\mathbf{w}_2)} =\innerprod[\mathcal{W}]{T\mathbf{v},a\mathbf{w}_1+\mathbf{w}_2} =\overline{a}\innerprod[\mathcal{W}]{T\mathbf{v},\mathbf{w}_1}+\innerprod[\mathcal{W}]{T\mathbf{v},\mathbf{w}_2} =\innerprod[\mathcal{V}]{\mathbf{v},aT^*\mathbf{w}_1+T^*\mathbf{w}_2} \nonumber
\end{align}
for every $\mathbf{v}\in\mathcal{V}$. Hence $T^*(a\mathbf{w}_1+\mathbf{w}_2)=aT^*\mathbf{w}_1+T^*\mathbf{w}_2$.\\
\hlt{Structural Commentary}: Each $\mathbf{w}$ gives a Riesz-represented functional on $\mathcal{V}$.\\
\hlt{Key Insight}: The adjoint is pointwise Riesz representation for $T$.
\end{proof}
\end{theorem}

\begin{figure}[H]
  \centering
  \begin{tikzpicture}[
  >=Stealth,
  every node/.style={font=\small},
  space/.style={draw=gray!45, rounded corners=3pt, minimum width=3.6cm, minimum height=2.4cm},
  vec/.style={->, blue!75!black, line width=0.85pt},
  adj/.style={->, teal!75!black, line width=0.85pt},
  map/.style={->, gray!70, line width=0.75pt}
]
  \node[space] (V) at (0,0) {};
  \node[anchor=north west] at (-1.6,1.05) {$\mathcal{V}$};
  \node[space] (W) at (5.2,0) {};
  \node[anchor=north west] at (3.6,1.05) {$\mathcal{W}$};

  \coordinate (Ov) at (-0.95,-0.65);
  \coordinate (v) at (0.55,0.55);
  \coordinate (Tsw) at (0.85,-0.25);
  \draw[vec] (Ov) -- (v) node[midway, above left] {$\mathbf{v}$};
  \draw[adj] (Ov) -- (Tsw) node[midway, below] {$T^*\mathbf{w}$};

  \coordinate (Ow) at (4.25,-0.65);
  \coordinate (Tv) at (5.85,0.55);
  \coordinate (w) at (5.95,-0.25);
  \draw[vec] (Ow) -- (Tv) node[midway, above left] {$T\mathbf{v}$};
  \draw[adj] (Ow) -- (w) node[midway, below] {$\mathbf{w}$};

  \draw[map] (1.9,0.45) -- (3.25,0.45) node[midway, above] {$T$};
  \draw[map] (3.25,-0.45) -- (1.9,-0.45) node[midway, below] {$T^*$};

  \node[draw=gray!45, rounded corners=3pt, align=center, inner sep=5pt] at (2.6,-1.85)
    {$\innerprod[\mathcal{W}]{T\mathbf{v},\mathbf{w}}
      =\innerprod[\mathcal{V}]{\mathbf{v},T^*\mathbf{w}}$};
\end{tikzpicture}

  \caption{The adjoint transfers $T$ between inner-product slots.}
  \label{fig:adjoint-inner-product-bridge}
\end{figure}

\begin{definition} \label{def:conjugate-transpose} \hlt{Conjugate Transpose}\\
For $A=[a_{ij}]\in\operatorname{Mat}(m,n,\mathbb{F})$, the \hlt{conjugate transpose} of $A$ is $A^*=[\overline{a_{ji}}]\in\operatorname{Mat}(n,m,\mathbb{F})$. Over $\mathbb{R}$, this is the ordinary transpose.
\end{definition}

\begin{theorem} \label{thm:matrix-of-adjoint-is-conjugate-transpose} \hlt{Matrix of the Adjoint}\\
Let $\mathcal{V}$ and $\mathcal{W}$ be finite-dimensional inner-product spaces. Let $\beta=(\mathbf{e}_1,\ldots,\mathbf{e}_n)$ be an orthonormal basis of $\mathcal{V}$ and let $\gamma=(\mathbf{f}_1,\ldots,\mathbf{f}_m)$ be an orthonormal basis of $\mathcal{W}$. If $T\in\LL(\mathcal{V},\mathcal{W})$, then $[T^*]_{\beta\leftarrow\gamma} =\bigl([T]_{\gamma\leftarrow\beta}\bigr)^*$.
\begin{proof}
Let $A=[T]_{\gamma\leftarrow\beta}$ and $B=[T^*]_{\beta\leftarrow\gamma}$.\\
By Theorem \ref{thm:matrix-entries-from-orthonormal-codomain}, $B_{i,j}=\innerprod[\mathcal{V}]{T^*\mathbf{f}_j,\mathbf{e}_i} =\overline{\innerprod[\mathcal{V}]{\mathbf{e}_i,T^*\mathbf{f}_j}} =\overline{\innerprod[\mathcal{W}]{T\mathbf{e}_i,\mathbf{f}_j}} =\overline{A_{j,i}}$. Thus $B=A^*$.\\
\hlt{Structural Commentary}: Entries swap by moving $T$ across the inner product.\\
\hlt{Key Insight}: In orthonormal coordinates, taking adjoints is exactly taking conjugate transposes.
\end{proof}
\end{theorem}

\begin{proposition} \label{prop:basic-adjoint-identities} \hlt{Basic Adjoint Identities}\\
For finite-dimensional inner-product spaces and composable linear maps:
\begin{enumerate}[label=\roman*.]
\setlength{\itemsep}{0pt}
\item $(aS+T)^*=\overline{a}S^*+T^*$ for $S,T\in\LL(\mathcal{V},\mathcal{W})$ and $a\in\mathbb{F}$;
\item $(SR)^*=R^*S^*$ for $R\in\LL(\mathcal{U},\mathcal{V})$ and $S\in\LL(\mathcal{V},\mathcal{W})$;
\item $(T^*)^*=T$;
\item $I_{\mathcal{V}}^*=I_{\mathcal{V}}$.
\end{enumerate}
\begin{proof}
Verify the defining identity; uniqueness gives equality.
\begin{enumerate}[label=(\roman*)]
\setlength{\itemsep}{0pt}
\item For $\mathbf{v}\in\mathcal{V}$, $\mathbf{w}\in\mathcal{W}$, $\innerprod{(aS+T)\mathbf{v},\mathbf{w}} =a\innerprod{S\mathbf{v},\mathbf{w}}+\innerprod{T\mathbf{v},\mathbf{w}} =\innerprod{\mathbf{v},\overline{a}S^*\mathbf{w}+T^*\mathbf{w}}$, so $(aS+T)^*=\overline{a}S^*+T^*$.
\item For $\mathbf{u}\in\mathcal{U}$ and $\mathbf{w}\in\mathcal{W}$, $\innerprod{SR\mathbf{u},\mathbf{w}} =\innerprod{R\mathbf{u},S^*\mathbf{w}} =\innerprod{\mathbf{u},R^*S^*\mathbf{w}}$, so $(SR)^*=R^*S^*$.
\item Conjugating $\innerprod{T\mathbf{v},\mathbf{w}}=\innerprod{\mathbf{v},T^*\mathbf{w}}$ gives $\innerprod{T^*\mathbf{w},\mathbf{v}}=\innerprod{\mathbf{w},T\mathbf{v}}$, so $T$ is the adjoint of $T^*$.
\item Since $\innerprod{I\mathbf{v},\mathbf{w}}=\innerprod{\mathbf{v},I\mathbf{w}}$, the identity map is its own adjoint.
\end{enumerate}
\hlt{Structural Commentary}: Put both sides in one inner-product slot, then use uniqueness.\\
\hlt{Key Insight}: Like matrix transposes, adjoints reverse products and conjugate scalars.
\end{proof}
\end{proposition}

\begin{proposition} \label{prop:adjoint-kernel-range-orthogonality} \hlt{Adjoint Kernel-Range Orthogonality}\\
Let $\mathcal{V}$ and $\mathcal{W}$ be finite-dimensional inner-product spaces and let $T\in\LL(\mathcal{V},\mathcal{W})$.\\
Then $\ker(T^*)=(\operatorname{range}(T))^{\perp}$, $\operatorname{range}(T^*)=(\ker(T))^{\perp}$.\\
Equivalently, the nullspace of $T^*$ is the orthogonal complement of the range of $T$.
\begin{proof}
For $\mathbf{w}\in\mathcal{W}$,
\begin{align}
\mathbf{w}\in\ker(T^*)
&\Longleftrightarrow \innerprod[\mathcal{V}]{\mathbf{v},T^*\mathbf{w}}=0\quad\text{for every }\mathbf{v}\in\mathcal{V} \nonumber\\
&\Longleftrightarrow \innerprod[\mathcal{W}]{T\mathbf{v},\mathbf{w}}=0\quad\text{for every }\mathbf{v}\in\mathcal{V} \nonumber\\
&\Longleftrightarrow \mathbf{w}\in(\operatorname{range}(T))^{\perp}.
\nonumber
\end{align}
Applying this identity to $T^*$ gives $\ker(T)=(\operatorname{range}(T^*))^{\perp}$.\\
Taking orthogonal complements and using Corollary \ref{cor:double-orthogonal-complement} yields $\operatorname{range}(T^*)=(\ker(T))^{\perp}$.\\
\hlt{Structural Commentary}: Adjoints trade ranges for kernels.\\
\hlt{Key Insight}: A vector is killed by $T^*$ exactly when it is orthogonal to everything $T$ can produce.
\end{proof}
\end{proposition}

\begin{definition} \label{def:self-adjoint-hermitian} \hlt{Self-Adjoint and Hermitian}\\
An operator $T\in\LL(\mathcal{V})$ on a finite-dimensional inner-product space is \hlt{self-adjoint} if $T^*=T$. A square matrix $A$ is \hlt{Hermitian} if $A^*=A$. In orthonormal coordinates, these are the same condition by Theorem \ref{thm:matrix-of-adjoint-is-conjugate-transpose}.
\end{definition}
